% HeavyBall replacement for the Richardson decoder sections.
% This fragment assumes theorem, proposition, proof, equation, and align are available.

\section{A one-head HeavyBall transformer decoder}
\label{sec:heavyball-decoder}

Let the encoder return an SPD matrix $\widehat A\succeq a_0 I$ and let
$P\succ0$ be an inverse preconditioner fixed during one recurrent forward
pass.  We initialize $z_{-1}=z_0=0$ and use the tied recurrence
\begin{equation}
 z_{\ell+1}
 =z_\ell+\alpha P(f-\widehat A z_\ell)
 +\beta(z_\ell-z_{\ell-1}),
 \qquad 0\leq \ell<L.
 \label{eq:hb-decoder}
\end{equation}
The two learned solver parameters $(\alpha,\beta)$ are shared over depth.

\begin{proposition}[HeavyBall from one Richardson attention block]
\label{prop:one-head-hb}
Suppose a tied attention block implements the Richardson map
\[
 \mathcal R_{\mathcal C}(z)
 =z+P_{\mathcal C}(f_{\mathcal C}-\widehat A_{\mathcal C}z)
\]
on one state channel.  After augmenting every state token by a delayed channel
$z_{\ell-1}$, one application of the same attention block followed by the
fixed linear channel map
\[
 (\mathcal R_{\mathcal C}(z_\ell),z_\ell,z_{\ell-1})
 \longmapsto
 \alpha\mathcal R_{\mathcal C}(z_\ell)
 +(1-\alpha+\beta)z_\ell-\beta z_{\ell-1}
\]
implements one step of \eqref{eq:hb-decoder}.  Thus HeavyBall adds no adaptive
inner product, reciprocal, or nonlinear arithmetic beyond what is already
required by the Richardson attention primitive.

For a strictly positive kernel, one normalized attention head is sufficient:
the same row-stochastic kernel matrix acts simultaneously on any number of
value channels, including the posterior-mean and posterior-variance
recursions.  Multiple heads are unnecessary unless distinct kernel operators
or mixture components must be advanced in parallel.
\end{proposition}

\begin{proof}
The identity
\[
 z_\ell+\alpha P(f-\widehat A z_\ell)
 +\beta(z_\ell-z_{\ell-1})
 =\alpha\mathcal R_{\mathcal C}(z_\ell)
 +(1-\alpha+\beta)z_\ell-\beta z_{\ell-1}
\]
gives the first claim.  Attention weights are shared across value coordinates,
so appending value channels changes neither the attention matrix nor the
number of heads.
\end{proof}

The last statement is specific to a positive normalized kernel.  Ordinary
softmax attention is row stochastic and therefore does not, in general,
represent the signed unnormalized weak-form moment $G^\top r$.  For the primal
weak-form system, Proposition~\ref{prop:one-head-hb} is invoked relative to an
exact linear/bilinear attention implementation of the Richardson map, or after
a positive-kernel reformulation.  We do not identify a softmax score
correlation with an exact matrix--vector product.

\section{Exact polynomial and finite-depth stability}
\label{sec:hb-polynomial}

Set
\[
 S=P^{1/2}\widehat A P^{1/2},
 \qquad \sigma(S)\subseteq[m,M],
 \qquad \kappa=M/m.
\]
If $e_\ell=z_\ell-\widehat A^{-1}f$ and
$\widetilde e_\ell=P^{-1/2}e_\ell$, then
\begin{equation}
 \widetilde e_{\ell+1}
 =\bigl((1+\beta)I-\alpha S\bigr)\widetilde e_\ell
 -\beta\widetilde e_{\ell-1}.
 \label{eq:hb-error-state}
\end{equation}
Define $p_{-1}(\lambda)=p_0(\lambda)=1$ and
\begin{equation}
 p_{\ell+1}(\lambda)
 =(1+\beta-\alpha\lambda)p_\ell(\lambda)
 -\beta p_{\ell-1}(\lambda).
 \label{eq:hb-polynomial}
\end{equation}
Then
\begin{equation}
 e_L=P^{1/2}p_L(S)P^{-1/2}e_0,
 \qquad
 z_L=\bigl[I-p_L(P\widehat A)\bigr]\widehat A^{-1}f.
 \label{eq:hb-linear-decoder}
\end{equation}
In particular, the finite-depth decoder remains linear in $f$.  This is the
essential difference from PCG: the PCG polynomial depends on the right-hand
side through adaptive Krylov coefficients.

\begin{theorem}[Optimal quadratic HeavyBall bound]
\label{thm:hb-finite-depth}
Choose
\begin{equation}
 \alpha_\star=\frac{4}{(\sqrt M+\sqrt m)^2},
 \qquad
 \beta_\star=
 \left(\frac{\sqrt M-\sqrt m}{\sqrt M+\sqrt m}\right)^2,
 \qquad
 q=\sqrt{\beta_\star}
 =\frac{\sqrt\kappa-1}{\sqrt\kappa+1}.
 \label{eq:hb-optimal-coefficients}
\end{equation}
For $z_{-1}=z_0=0$ and every $L\geq0$,
\begin{equation}
 \|e_L\|_{P^{-1}}
 \leq C_L(q)q^L\|e_0\|_{P^{-1}},
 \qquad
 C_L(q)=1+L(1+q).
 \label{eq:hb-transient-bound}
\end{equation}
More generally, constant coefficients are strictly stable on $[m,M]$ when
\begin{equation}
 0\leq\beta<1,
 \qquad
 0<\alpha M<2(1+\beta).
 \label{eq:hb-stability-region}
\end{equation}
\end{theorem}

\begin{proof}
For \eqref{eq:hb-optimal-coefficients}, write
$1+\beta_\star-\alpha_\star\lambda=2qt$ with $t\in[-1,1]$.
After setting $s_\ell=p_\ell/q^\ell$, recurrence
\eqref{eq:hb-polynomial} becomes
$s_{\ell+1}=2t s_\ell-s_{\ell-1}$.  Hence
\[
 p_L(\lambda)=q^L\bigl(U_L(t)-qU_{L-1}(t)\bigr),
\]
where $U_L$ is the Chebyshev polynomial of the second kind.  Since
$|U_L(t)|\leq L+1$ on $[-1,1]$,
\[
 \max_{\lambda\in[m,M]}|p_L(\lambda)|
 \leq q^L\{L+1+qL\}=C_L(q)q^L.
\]
Functional calculus in \eqref{eq:hb-error-state} proves
\eqref{eq:hb-transient-bound}.  The strict stability region follows from the
Jury conditions for
$r^2-(1+\beta-\alpha\lambda)r+\beta$ uniformly over
$\lambda\in[m,M]$.
\end{proof}

The factor $C_L(q)$ is necessary in a uniform finite-depth statement because
the characteristic roots coalesce at the endpoints $m$ and $M$.  A bare
$q^L$ bound with constant one is false.  Numerical plots should use the exact
polynomial \eqref{eq:hb-polynomial}; \eqref{eq:hb-transient-bound} is a safe
closed-form envelope.

\section{Operator-supervised encoder--decoder transfer}
\label{sec:hb-transfer}

\begin{theorem}[HeavyBall encoder--decoder transfer]
\label{thm:hb-transfer}
Let $A,\widehat A\succeq a_0I$, let $f\sim\mathcal N(0,\Sigma)$ be conditionally
independent of the encoder sample, and suppose
$\kappa(P)\leq\overline\kappa_P$.  Let $z_L$ be defined by
\eqref{eq:hb-decoder} with the optimal coefficients for a uniform spectral
interval $[m,M]$ containing
$\sigma(P^{1/2}\widehat A P^{1/2})$.  Define
\[
 \mathcal R_A=\frac1D\mathbb E\|\widehat A-A\|_F^2,
 \qquad
 \mathcal R_{u,L}^{\rm HB}=\mathbb E\|z_L-A^{-1}f\|_2^2.
\]
Then
\begin{equation}
 \boxed{
 \mathcal R_{u,L}^{\rm HB}
 \leq
 \frac{2\overline\kappa_P\,\operatorname{tr}(\Sigma)}{a_0^2}
 C_L(q)^2q^{2L}
 +\frac{2D\operatorname{tr}(\Sigma)}{a_0^4}\mathcal R_A.}
 \label{eq:hb-transfer}
\end{equation}
Conditionally on $(A,\widehat A,P)$, the exact identity is
\begin{equation}
 \mathbb E_f\|z_L-A^{-1}f\|_2^2
 =\operatorname{tr}\!\left(\Sigma T_L^\top T_L\right),
 \quad
 T_L=\bigl[I-p_L(P\widehat A)\bigr]\widehat A^{-1}-A^{-1}.
 \label{eq:hb-exact-trace}
\end{equation}
\end{theorem}

\begin{proof}
Equation \eqref{eq:hb-linear-decoder} gives
\[
 \|z_L-\widehat A^{-1}f\|_2
 \leq
 \sqrt{\kappa(P)}\,C_L(q)q^L
 \|\widehat A^{-1}f\|_2
 \leq
 \frac{\sqrt{\overline\kappa_P}C_L(q)q^L}{a_0}\|f\|_2.
\]
The resolvent identity gives
\[
 \|\widehat A^{-1}f-A^{-1}f\|_2
 \leq a_0^{-2}\|\widehat A-A\|_{\rm op}\|f\|_2.
\]
Use $(a+b)^2\leq2a^2+2b^2$,
$\mathbb E\|f\|_2^2=\operatorname{tr}(\Sigma)$,
$\|\widehat A-A\|_{\rm op}^2\leq\|\widehat A-A\|_F^2$,
and conditional independence.  The trace identity follows because $T_L$ is
linear in $f$.
\end{proof}

Compared with the Richardson transfer theorem, the encoder term is unchanged.
Only the decoder polynomial changes.  Compared with PCG, HeavyBall preserves
the exact conditional Gaussian trace formula because its coefficients are
fixed independently of $f$.

\section{Spectral and replica-ready form}
\label{sec:hb-replica}

The second-order recurrence is first order on the augmented state:
\begin{equation}
 \begin{bmatrix}e_{\ell+1}\\e_\ell\end{bmatrix}
 =
 \underbrace{\begin{bmatrix}
 (1+\beta)I-\alpha P\widehat A&-\beta I\\
 I&0
 \end{bmatrix}}_{\mathbb T_{\rm HB}}
 \begin{bmatrix}e_\ell\\e_{\ell-1}\end{bmatrix}.
 \label{eq:hb-augmented-state}
\end{equation}
Thus the decoder order parameter is the spectral filter
$p_L(\lambda)^2$, not a new adaptive Krylov overlap at every layer.  In a
commuting model with isotropic query covariance
$\Sigma=\sigma_f^2I$ and joint eigenvalues
$(\widehat a_i,p_i)$, the solver part of the exact risk is
\begin{equation}
 \mathcal R_{\rm solve,L}^{\rm HB}
 =\sigma_f^2\sum_{i=1}^D
 \frac{p_L(p_i\widehat a_i)^2}{\widehat a_i^2}.
 \label{eq:hb-spectral-risk}
\end{equation}
If the empirical joint spectral measure converges to $\nu$, then the
normalized risk converges formally to
\begin{equation}
 \frac1D\mathcal R_{\rm solve,L}^{\rm HB}
 \longrightarrow
 \sigma_f^2\int \frac{p_L(p\widehat a)^2}{\widehat a^2}
 \,d\nu(\widehat a,p).
 \label{eq:hb-rmt-integral}
\end{equation}
This is the direct replacement of the Richardson filter
$(1-p\widehat a)^{2L}$ in the random-matrix and replica calculations.  A
replica calculation is still required for the learned encoder and learned
preconditioner order parameters; no replica machinery is needed merely to
describe the fixed HeavyBall decoder polynomial.
